\begin{textsample}{POS Dim 1 – 2017\_c5 – Score 15.00 – 2017\_c5/enade\_2017\_art\_vis\_8100b5f9fc.txt}  \label{ex:f1_pos_060}
CAPÍTULO 5 DISTRIBUIÇÃO DOS CONCEITOS 5.1 \textbf{PANORAMA} NACIONAL DA DISTRIBUIÇÃO DOS CONCEITOS O Conceito Enade é um indicador de qualidade que avalia os cursos por \textbf{intermédio} dos desempenhos dos estudantes a partir dos resultados do Enade.
Os valores possíveis vão de 1 ( pior situação ) a 5 ( melhor situação ).
Os cursos com apenas um ou sem nenhum concluinte participante não obtêm o Conceito Enade, ficando Sem Conceito ( SC ).
A Tabela 5.1 apresenta a quantidade e distribuição de cursos de Artes Visuais ( Licenciatura ) participantes do Enade/2017, por faixa de conceito e Grande Região.
Enfatiza -se, mais uma vez, que a diferença entre os cursos \textbf{tabulados} neste capítulo e no capítulo 2 corresponde aos cursos Sem Conceito ( SC ), ou seja, àqueles, em princípio, sem alunos concluintes que houvessem participado da prova ou àqueles que tivessem tão somente um aluno concluinte18.
Observando -se os dados da Tabela 5.1, nota -se que, dos 102 cursos participantes, 33 ( 32,4 \% ) classificaram -se com conceito 3, o valor modal.
Esse foi também o conceito modal na região Sul, com 11 cursos ( 36,7 \% ).
Na região Nordeste, houve uma distribuição \textbf{bimodal}, já que cinco cursos receberam os conceitos 2 e 3 ( 35,7 \% para cada um ).
Na região Norte, os conceitos 1, 2 e 3 foram igualmente atribuídos a dois cursos cada, correspondendo a 28,6 \% para cada um.
Na região Sudeste, 15 cursos ( 32,6 \% ) receberam o conceito modal 2.
Já na região Centro-Oeste, a moda foi o conceito 4, atribuído a dois cursos ( 40,0 \% ).
O conceito 2 foi o segundo mais frequente em nível nacional ( 25,5 \%, correspondendo a 26 cursos ), e o conceito 4, o terceiro ( 16,7 \%, correspondendo a 17 cursos ).
Houve, ainda, 14 cursos ( 13,7 \% ) que receberam conceito 1, e oito cursos ( 7,8 \% ) que receberam conceito 5.
Dos 102 cursos de Artes Visuais ( Licenciatura ), quatro ( 3,9 \% ) ficaram Sem Conceito ( SC ). 18 Esses cursos com somente um concluinte não foram considerados no capítulo 2.
A região Norte participou com sete cursos ou 6,9 \% do total nacional.
Desses, seis cursos foram distribuídos entre os conceitos 1, 2 e 3, ou seja, cada um desses conceitos foi representado por dois cursos ( 28,6 \%, cada um ).
O outro curso \textbf{restante} foi alocado ao conceito 4 ( 14,3 \% ).
Nenhum curso ficou Sem Conceito ( SC ), e nenhum curso recebeu conceito 5 nessa região.
A região Nordeste participou com 14 cursos ou 13,7 \% do total nacional.
Nessa região, houve distribuição \textbf{bimodal} entre os conceitos 2 e 3, com cinco cursos, cada um ( 35,7 \% para cada um ).
Três cursos ( 21,4 \% em termos regionais ) obtiveram conceito 4, e um curso obteve conceito 1.
Nenhum curso ficou Sem Conceito ( SC ), assim como nenhum recebeu conceito 5.
Dos 46 cursos participantes da região Sudeste ( 45,1 \% do total nacional ), o conceito 2 foi o conceito modal, como já anteriormente assinalado, tendo sido obtido por 15 cursos ( 32,6 \% ).
O conceito 3 foi atribuído a 14 cursos ( 30,4 \% ), e o conceito 1, a seis cursos ( 13,0 \% ).
Receberam os conceitos 4 e 5, respectivamente, quatro cursos ( 8,7 \% ) e cinco cursos ( 10,9 \% ).
Nessa região, dois cursos ficaram Sem Conceito ( SC ).
Os 30 cursos da região Sul corresponderam a 29,4 \% do total nacional.
Como já \textbf{indicado}, houve predominância do conceito 3, atribuído a 36,7 \% dos cursos ( conceito modal ), ou seja, atribuído a 11 dos 30 cursos participantes na região Sul.
Os conceitos 1 e 2 foram atribuídos a quatro cursos, cada um ( 13,3 \% ), o conceito 4, a sete cursos ( 23,3 \% ), e o conceito 5, a dois cursos ( 6,7 \% ).
Dois cursos ficaram Sem Conceito ( SC ) na região Sul.
Já dos cinco cursos participantes na região Centro-Oeste ( 4,9 \% do total nacional ), dois receberam conceito 4, o conceito modal ( 40,0 \% ), como já fora destacado.
Os conceitos \textbf{restantes}, 1, 3 e 5, foram atribuídos a um curso, cada.
Nessa região, nenhum curso ficou Sem Conceito ( SC ), assim como nenhum recebeu conceito 2.
No Gráfico 5.1, é apresentada a distribuição \textbf{cumulativa} do Conceito Enade segundo a Grande Região.
Em um gráfico de distribuição \textbf{cumulativa}, em cada valor de \textbf{abscissa} ( eixo x ), o valor apresentado no eixo das \textbf{ordenadas} ( eixo y ) é o valor acumulado da \textbf{variável}.
Por exemplo, a região Norte ( linha azul ) apresenta um pouco mais de 85 \% dos cursos com conceito 3 ou abaixo deste valor.
A região Centro-Oeste ( linha preta ), por outro lado, apresenta 40 \% dos cursos na mesma situação : com conceito 3 ou abaixo.
Nas regiões Norte ( linha azul ) e Nordeste ( linha verde ), todos os cursos têm conceitos 4 ou abaixo disso, já que a \textbf{poligonal} \textbf{cumulativa} alcança 100 \% para o conceito 4.
Em linhas gerais, num gráfico de distribuição \textbf{cumulativa}, \textbf{poligonais} mais à esquerda representam Grandes Regiões com pior distribuição, e \textbf{poligonais} mais à direita, Grandes Regiões com melhores conceitos.
A região Norte apresenta os cursos com conceitos mais baixos, ao passo que a região Centro-Oeste ( linha preta ) apresenta a distribuição com valores maiores.
Cursos em IES situadas nas regiões Sul ( linha \textbf{roxa} ) e Sudeste ( linha vermelha ) apresentam uma situação intermediária. 5.2 CONCEITOS POR CATEGORIA ADMINISTRATIVA, POR MODALIDADE DE ENSINO E GRANDE REGIÃO Na Tabela 5.2, é apresentada a distribuição dos cursos de Artes Visuais ( Licenciatura ) participantes do Enade/2017, por Categoria Administrativa e Modalidade de Ensino, de acordo com os conceitos por eles alcançados, segundo as Grandes Regiões brasileiras.
Dos 102 cursos participantes, 55 ( 53,9 \% ) eram \textbf{ministrados} em Instituições \textbf{Públicas}, e 47 ( 46,1 \% ), em Instituições Privadas.
Quanto à Modalidade de Ensino, 86,3 \% dos cursos eram Presenciais, e 13,7 \%, a Distância.
De acordo com as informações da Tabela 5.2, em termos nacionais, entre Instituições \textbf{Públicas} e Privadas, observa -se que, dos oito cursos avaliados com conceito 5, sete eram oferecidos em IES Públicas e um, em IES Privadas.
Dos 55 cursos participantes de IES Públicas, o conceito 3 foi o valor modal, atribuído a 15 cursos ( 27,3 \% ).
Entre os demais cursos de Instituições \textbf{Públicas} participantes, oito cursos ( 14,5 \% ) obtiveram conceito 1, dez cursos ( 18,2 \% ) obtiveram conceito 2, 14 cursos ( 25,5 \% ), conceito 4 e, como já comentado, sete cursos ( 12,7 \% ), conceito 5.
Nessa categoria, um curso ficou Sem Conceito ( SC ).
Na rede privada, o conceito modal foi 3, atribuído a 18 cursos, o correspondente a 38,3 \% dos 47 cursos da categoria.
Entre os demais cursos participantes, seis cursos ( 12,8 \% ) foram avaliados com conceito 1, 16 cursos ( 34,0 \% ) receberam conceito 2.
O conceito 4 foi atribuído a três cursos ( 6,4 \% ), e o conceito 5, a um curso ( 2,1 \% ).
Três cursos ( 6,4 \% ) ficaram Sem Conceito ( SC ).
Dos cursos participantes da modalidade Presencial, o conceito 3 foi o valor modal, atribuído a 35,2 \% dos cursos.
Entre os demais cursos nessa Modalidade de Ensino, 10,2 \% obtiveram conceito 1, 21,6 \% obtiveram conceito 2, 19,3 \% receberam conceito 4, e 9,1 \% dos cursos, conceito 5.
Nessa modalidade, quatro cursos ficaram Sem Conceito ( SC ).
Entre os cursos a Distância, o conceito modal foi 2, atribuído a 50,0 \% dos cursos da modalidade.
Dentre os demais cursos participantes, 35,7 \% receberam conceito 1, e 14,3 \%, conceito 2.
Nessa modalidade, nenhum curso ficou Sem Conceito ( SC ) ou recebeu os conceitos 4 ou 5.
Na análise por região, observa -se que, na região Norte, as Instituições Privadas participaram com apenas um curso ( 14,3 \% do total regional ), o qual recebeu conceito 3.
As Instituições Públicas participaram com seis cursos na região Norte ( 85,7 \% do total regional ).
Os conceitos modais para as Instituições \textbf{Públicas} ficaram distribuídos entre os conceitos 1 e 2, tendo sido atribuídos a dois cursos, cada um ( 33,3 \% para cada um ).
Os dois cursos \textbf{restantes} também foram distribuídos entre os conceitos 3 e 4, com um curso cada.
Nenhum curso ficou Sem Conceito ( SC ), assim como nenhum recebeu conceito 5 nessa combinação de categoria e região.
Todos os cursos oferecidos na região Norte eram Presenciais.
Na região Nordeste, a rede privada \textbf{contou} com apenas um curso participante, o equivalente a 7,1 \% do total da região.
Esse curso foi avaliado com conceito 2.
As Instituições \textbf{Públicas} da região Nordeste participaram com 13 cursos ( 92,9 \% do total da região ).
Desses, cinco obtiveram conceito 3, o valor modal.
Um curso recebeu conceito 1, quatro cursos receberam conceito 2, e três cursos receberam conceito 4.
Nenhum curso ficou Sem Conceito ( SC ) ou com conceito 5 nessa região.
A maioria dos cursos na região foram Presenciais : 85,7 \%.
Para esses, a moda foi o conceito 3 com 41,7 \% cursos dessa modalidade.
Os demais cursos Presenciais da região foram alocados aos conceitos 2 ( 33,3 \% ), e 4, 25,0 \% dos cursos.
Nenhum curso Presencial ficou Sem Conceito ( SC ), assim como nenhum recebeu os conceitos 1 ou 5.
Os cursos a Distância foram distribuídos igualmente entre os conceitos 1 e 2, ou seja, 50,0 \% dos cursos para cada um.
Nenhum curso a Distância da região foi alocado aos demais conceitos ou ficou Sem Conceito ( SC ).
Na região Sudeste, a proporção de cursos da rede privada, 65,2 \%, foi mais elevada do que a registrada nas demais regiões brasileiras, correspondendo a 30 dos 46 cursos participantes.
Nessa Categoria e Região, houve distribuição \textbf{bimodal} entre os conceitos 2 e 3, atribuídos a 11 cursos, cada um.
Dois cursos ficaram Sem Conceito ( SC ).
Os demais foram avaliados com conceito 1 ( quatro cursos ), 4 ( um curso ), e conceito 5 ( um curso também ).
Entre os 16 cursos oferecidos por Instituições \textbf{Públicas} na região Sudeste, ocorreu distribuição \textbf{bimodal} entre os conceitos 2 e 5, atribuídos a quatro cursos, cada um.
Os demais cursos receberam os conceitos 3 e 4 ( ambos com três cursos ), e conceito 1 ( dois cursos ).
Na região Sudeste, nessa categoria, nenhum curso ficou Sem Conceito ( SC ).
Também na região Sudeste, a maioria dos cursos foram Presenciais : 80,4 \%.
Para esses, a moda foi o conceito 3 com 32,4 \%.
Os demais cursos Presenciais da região foram alocados aos conceitos 1 ( 10,8 \% ), 2 ( 27,0 \% ), 4 ( 10,8 \% ) e 5 ( 13,5 \% ).
Dois cursos Presenciais ficaram Sem Conceito ( SC ).
Para os cursos a Distância, a moda foi o conceito 2, atribuído a 55,6 \% dos cursos.
Os demais cursos a Distância da região foram alocados aos conceitos 1 ( 22,2 \% ), e 3 ( 22,2 \% ).
Nenhum curso a Distância da região foi alocado aos conceitos 4 ou 5, e nenhum ficou Sem Conceito ( SC ).
As Instituições Privadas concentraram 15 dos 30 cursos participantes da região Sul, 50,0 \% do total regional.
Desses, seis obtiveram conceito 3, o conceito modal.
Nessa combinação de Categoria Administrativa e Grande Região, os conceitos 1 e 4 foram atribuídos a dois cursos, cada um.
Quatro cursos receberam conceito 2, e um curso ficou Sem Conceito ( SC ).
O conceito 5 não foi atribuído a nenhum curso nessa região.
As Instituições \textbf{Públicas} na região Sul participaram também com 15 cursos, 50,0 \% dos cursos da região.
O conceito modal foi distribuído entre os conceitos 3 e 4, que foram atribuídos a cinco cursos, cada um.
Os conceitos 1 e 5 foram também atribuídos a dois cursos, cada um.
Nenhum curso da região Sul oferecido por IES Pública recebeu conceito 2, e um curso ficou Sem Conceito ( SC ).
Também na região Sul, a maioria dos cursos foram Presenciais, correspondendo a 93,3 \%.
Para os cursos Presenciais, a moda foi o conceito 3, com 39,3 \% dos cursos.
Os demais cursos Presenciais da região foram alocados aos conceitos 4 ( 25,0 \% ) e aos conceitos 1 e 2 ( 10,7 \%, cada um ).
Dois cursos Presenciais foram alocados ao conceito 5 e os outros dois cursos ficaram Sem Conceito ( SC ).
Dos cursos a Distância, metade deles recebeu conceito 1, e a outra metade recebeu conceito 2.
Na região Centro-Oeste, todos os cinco cursos participantes eram de Instituições \textbf{Públicas} ( 100,0 \% em termos regionais ).
Desses, dois receberam conceito 4, o conceito modal.
Entre os três \textbf{restantes}, cada um recebeu os conceitos 1, 3 e 5.
Nenhum curso ficou Sem Conceito ( SC ) ou com conceito 2.
Não houve cursos oferecidos por Instituições \textbf{Públicas} na região Centro-Oeste.
Nesta região, 80,0 \% dos cursos foram Presenciais e 20,0 \% a Distância.
O conceito modal para os cursos Presenciais foi 4, atribuído a dois dos quatro cursos nessa modalidade.
Os demais cursos receberam conceitos 3 e 5.
Os conceitos 1 e 2 não foram atribuídos a nenhum curso, assim como nenhum curso ficou Sem Conceito ( SC ).
Ao único curso a Distância foi atribuído o conceito 1.
O Gráfico 5.2 apresenta a distribuição do Conceito Enade segundo a Categoria Administrativa da IES.
Os cursos em IES Públicas ( linha azul ) apresentam uma \textbf{poligonal} mais à direita do que os em IES Privadas ( linha verde ) e, consequentemente, uma distribuição de Conceitos Enade com valores maiores.
Nota -se que cerca de 60 \% dos cursos em IES Públicas apresentam conceito 3 ou abaixo disso, enquanto os cursos em IES Privadas nesta mesma condição são cerca de 90 \%.
Na combinação de Categoria Administrativa e Grande Região, os conceitos 1 e 4 foram atribuídos a dois cursos, cada um.
Quatro cursos receberam conceito 2, e um curso ficou Sem Conceito ( SC ).
O conceito 5 não foi atribuído a nenhum curso nessa região.
As Instituições \textbf{Públicas} na região Sul participaram também com 15 cursos, 50,0 \% dos cursos da região.
O conceito modal foi distribuído entre os conceitos 3 e 4, que foram atribuídos a cinco cursos, cada um.
Os conceitos 1 e 5 foram também atribuídos a dois cursos, cada um.
Nenhum curso da região Sul oferecido por IES Pública recebeu conceito 2, e um curso ficou Sem Conceito ( SC ).
Também na região Sul, a maioria dos cursos foram Presenciais, correspondendo a 93,3 \%.
Para os cursos Presenciais, a moda foi o conceito 3, com 39,3 \% dos cursos.
Os demais cursos Presenciais da região foram alocados aos conceitos 4 ( 25,0 \% ) e aos conceitos 1 e 2 ( 10,7 \%, cada um ).
Dois cursos Presenciais foram alocados ao conceito 5 e os outros dois cursos ficaram Sem Conceito ( SC ).
Dos cursos a Distância, metade deles recebeu conceito 1, e a outra metade recebeu conceito 2.
Na região Centro-Oeste, todos os cinco cursos participantes eram de Instituições \textbf{Públicas} ( 100,0 \% em termos regionais ).
Desses, dois receberam conceito 4, o conceito modal.
Entre os três \textbf{restantes}, cada um recebeu os conceitos 1, 3 e 5.
Nenhum curso ficou Sem Conceito ( SC ) ou com conceito 2.
Não houve cursos oferecidos por Instituições \textbf{Públicas} na região Centro-Oeste.
Nesta região, 80,0 \% dos cursos foram Presenciais e 20,0 \% a Distância.
O conceito modal para os cursos Presenciais foi 4, atribuído a dois dos quatro cursos nessa modalidade.
Os demais cursos receberam conceitos 3 e 5.
Os conceitos 1 e 2 não foram atribuídos a nenhum curso, assim como nenhum curso ficou Sem Conceito ( SC ).
Ao único curso a Distância foi atribuído o conceito 1.
O Gráfico 5.2 apresenta a distribuição do Conceito Enade segundo a Categoria Administrativa da IES.
Os cursos em IES Públicas ( linha azul ) apresentam uma \textbf{poligonal} mais à direita do que os em IES Privadas ( linha verde ) e, consequentemente, uma distribuição de Conceitos Enade com valores maiores.
Nota -se que cerca de 60 \% dos cursos em IES Públicas apresentam conceito 3 ou abaixo disso, enquanto os cursos em IES Privadas nesta mesma condição são cerca de 90 \%. 5.3 CONCEITOS POR ORGANIZAÇÃO ACADÊMICA E GRANDE REGIÃO Na Tabela 5.3, encontra -se a distribuição dos conceitos atribuídos aos cursos participantes do Enade/2017 na Área de Artes Visuais ( Licenciatura ), por Organização Acadêmica, segundo as Grandes Regiões brasileiras.
Dos 102 cursos de Artes Visuais ( Licenciatura ) participantes, 75 eram oferecidos em Universidades, dez em Centros Universitários, 14 em Faculdades e três em CEFET/IFET.
Essa distribuição corresponde a, respectivamente, 73,5 \%, 9,8 \%, 13,7 \% e 2,9 dos cursos.
De acordo com os dados apresentados, de todos os oito cursos avaliados com conceito 5, sete eram vinculados a Universidades.
Os cursos vinculados a esse tipo de Organização Acadêmica tiveram o conceito 3 como conceito modal, atribuído a 25 cursos ( 33,3 \% ).
Os demais cursos vinculados a Universidades receberam os conceitos 1 ( oito cursos ), 2 ( 19 cursos ), 4 ( 14 cursos ) e 5 ( sete cursos, como já mencionado ).
Ficaram Sem Conceito ( SC ) dois cursos.
Dentre os cursos em Centros Universitários, houve distribuição \textbf{bimodal} entre os conceitos 2 e 3, atribuídos a três cursos, cada um ( 30,0 \%, cada ).
Um curso vinculado a esse tipo de Organização Acadêmica ficou Sem Conceito ( SC ), e os outros cursos receberam os conceitos 1 ( dois cursos ) e 4 ( um curso ).
Nenhum curso recebeu conceito 5.
Dos 14 cursos mantidos por Faculdades, quatro receberam conceito 2, e outros quatro receberam conceito 3, em uma distribuição \textbf{bimodal}.
O segundo conceito mais frequentemente atribuído aos cursos oferecidos por essa categoria de IES foi o conceito 1, recebido por três cursos.
Na sequência, um curso obteve conceito 4, e um curso também, conceito 5.
Um curso ficou Sem Conceito ( SC ).
Nos três cursos mantidos por CEFET/IFET, os conceitos 1, 3 e 4 foram atribuídos, cada um, a um curso.
Sendo assim, nenhum curso ficou Sem Conceito ( SC ) ou recebeu os conceitos 2 ou 5.
Considerando -se separadamente as regiões brasileiras, verifica -se que, na região Norte, as Universidades participaram com todos os sete cursos oferecidos.
Os conceitos 1, 2 e 3 foram atribuídos a dois cursos, cada um.
O outro curso \textbf{restante} obteve conceito 4.
Nenhum curso ficou Sem Conceito ( SC ) ou recebeu conceito 5 nessa combinação de Categoria e Região.
Não houve cursos representantes nos Centros Universitários, nas Faculdades e nos CEFET/IFET da região Norte.
Na região Nordeste, as Universidades participaram com 11 dos 14 cursos da Área de Artes Visuais ( Licenciatura ) oferecidos na região, dos quais nenhum ficou Sem Conceito ( SC ).
Nota -se que tais cursos apresentaram uma distribuição com moda distribuída em dois conceitos contíguos, 2 e 3, atribuídos, cada qual, a quatro cursos.
Obtiveram os conceitos 1 e 4, respectivamente, um e dois cursos.
Nenhum curso oferecido por Universidades no Nordeste recebeu conceito 5.
Os Centros Universitários não tiveram cursos participantes na região Nordeste.
As Faculdades foram representadas por apenas um curso, que obteve conceito 2.
Os CEFET/IFET participaram com dois cursos, dos quais um recebeu conceito 3, e o outro recebeu conceito 4.
Na região Sudeste, as Universidades concentraram 28 dos 46 cursos de Artes Visuais ( Licenciatura ) da região.
Entre os cursos em Universidades, nessa região, o conceito modal foi 2, atribuído a nove cursos.
Os demais cursos receberam os conceitos 1 ( três cursos ), 3 ( oito cursos ), 4 ( três cursos ) e 5 ( quatro cursos ).
Um curso ficou Sem Conceito ( SC ).
Os Centros Universitários participaram com oito cursos na região Sudeste.
Os conceitos 2 e 3 foram distribuídos, cada um deles, por três cursos.
O conceito 1 foi atribuído a um curso.
Um curso ficou Sem Conceito ( SC ) nessa combinação de Organização Acadêmica e região.
Nenhum curso foi atribuído aos conceitos 4 ou 5.
Já as Faculdades foram representadas por dez cursos na região Sudeste, dos quais nenhum ficou Sem Conceito ( SC ).
Aos demais cursos foram atribuídos os conceitos 1 ( dois cursos ), 2 e 3, os valores modais ( três cursos para cada um ), 4 e 5 ( um curso para cada um ).
Os CEFET/IFET não participaram com nenhum curso.
Dos 30 cursos da região Sul, 24 eram de Universidades, para os quais o conceito modal foi 3, atribuído a dez cursos.
Nesse tipo de organização, um curso ficou Sem Conceito ( SC ).
Um curso recebeu conceito 1, quatro cursos receberam conceito 2, seis cursos, conceito 4, e dois cursos, conceito 5.
Dos dois cursos participantes de Centros Universitários da região Sul, nenhum ficou Sem Conceito ( SC ).
Um curso recebeu conceito 1 e o outro curso representante obteve conceito 4.
Quanto aos três cursos vinculados a Faculdades na região Sul, um ficou Sem Conceito ( SC ).
Entre os demais, um curso recebeu 1, e outro recebeu conceito 3.
Nessa região, apenas um curso foi avaliado em CEFET/IFET, o qual obteve conceito 1.
Na região Centro-Oeste, todos os cinco cursos participantes eram de Universidades.
Nesse tipo de organização, nenhum curso ficou Sem Conceito ( SC ), assim como nenhum recebeu conceito 2.
Os conceitos 1, 3 e 5 foram atribuídos a um curso cada.
Os outros dois cursos obtiveram conceito 4, conceito modal.
Os Centros Universitários, Faculdades e CEFET/IFET da região Centro-Oeste não tiveram cursos representantes.
No Gráfico 5.3, é apresentada a distribuição \textbf{cumulativa} do Conceito Enade segundo a Organização Acadêmica da IES.
Os cursos em CEFET/IFET ( linha \textbf{roxa} ) e em Universidades ( linha azul ) apresentam \textbf{poligonais} mais à direita do que as demais, \textbf{denotando} uma melhor distribuição de Conceitos Enade.
No outro extremo, encontram -se os cursos oferecidos em Centros Universitários ( linha verde ) e Faculdades ( linha vermelha ) com as \textbf{poligonais} mais à esquerda e, consequentemente, uma distribuição de Conceitos Enade com menores valores.
Cursos em Centros Universitários e em CEFET/IFET ( linha \textbf{roxa} ) não obtiveram conceito 5, pois as \textbf{poligonais} \textbf{cumulativas} alcançam 100 \% no conceito 4.
Os cursos em CEFET/IFET ( linha \textbf{roxa} ) também não obtiveram conceito 2, pois a \textbf{poligonal} se comporta como uma linha horizontal entre os conceitos 1 e 2.

% matched lemmas: abscissa, bimodal, contar, cumulativo, denotar, indicado, intermédio, ministrar, ordenado, panorama, poligonal, pública, restante, roxa, tabular, variável
\end{textsample}
